\documentclass{sdjumpracticereport}
\begin{document}
\begin{titlepage}
\centering
\vspace*{\dimexpr7mm-12bp\relax}
\CoverLogo
\vspace{10mm}
\CoverTitle{赴测试实践中心\\专业实践总结报告}
\vspace{11mm}
\begin{coverfields}
\CoverField{姓\quad 名}{测试研究生甲}
\CoverField{学\quad 号}{local:practice-report}
\CoverField{校内导师}{测试导师甲}
\CoverField{企业导师}{测试导师乙}
\CoverField{研究方向}{模拟设备状态监测}
\CoverField{类别（领域）}{电子信息（控制工程）}
\CoverField{学\quad 院}{测试工程学院}
\CoverField{实践时间}{2025年11月至2026年11月}
\end{coverfields}
\end{titlepage}

% 采用原件的六项目录，正文标题由下列各节直接填写。
{\centering\sdjufont{18}{24}目录\par}
\bigskip
\begin{enumerate}
\renewcommand\labelenumi{\chinese{enumi}、}
\setlength\itemsep{11bp}
\sdjufont{14}{28}\bfseries
\item 实践目的
\item 实践计划
\item 实践单位简介
\item 实践内容
\item 收获与体会
\item 实践评价
\end{enumerate}

\clearpage
\FormHeading{一、实践目的}
在模拟环境中练习台架数据记录和复核流程。
\FormHeading{二、实践计划}
先完成安全培训与测点核验，再分阶段采集、归档和复核。
\FormHeading{三、实践单位简介}
测试实践中心为虚构单位，示例中用于组织模拟台架实践。
\FormHeading{四、实践内容}
在三类模拟工况下记录设备信号，并整理实验配置和数据字典。
\FormHeading{五、收获与体会}
通过记录配置和操作时间，复核人员可以定位重复试验之间的差异。
\FormHeading{六、实践评价}
完成虚构实践计划所列任务，提交可复核的模拟记录。
\end{document}
